En este ejemplo en Python probamos cómo funciona UDP en localhost usando dos hilos: uno para el servidor y otro para el cliente.
El servidor abre un socket de tipo SOCK\_DGRAM en el puerto 5002 y queda esperando con recvfrom, sin hacer handshake ni aceptar una conexión previa.
El cliente espera 0.3 segundos solo para asegurarse de que el servidor ya hizo bind y luego manda el paquete directo con sendto.
Cuando el servidor recibe el datagrama, lee el mensaje junto con la dirección de origen y devuelve la respuesta por la misma vía.
Con esto se ve lo principal de UDP: se ahorra el handshake y las retransmisiones, pero a cambio no garantiza que el paquete llegue ni que llegue rápido.